\section{Introduction}

\begin{frame}
  \frametitle{Beamer の特徴：数式が書ける}
  
  基本的に，\LaTeX のメリット{\footnotesize （とデメリット）}がそのまま当てはまる
  \begin{itemize}
    \okitem \emph{数式を書くのが GUI 系のアプリより遥かに楽}
    \okitem \LaTeX の文書との連携が楽
    \ngitem 図を作るのはめちゃめちゃ大変\footcite{tikz-nakata}，コンパイルが重い
  \end{itemize}

  \begin{block}{However, }
    ゼミで一番ありがちな指摘は，「\Emph{定義は?}」
    \begin{itemize}
      \thusitem \emph{数式が書けるというのは大きい}
    \end{itemize}
  \end{block}
\end{frame}

\begin{frame}
  \frametitle{使い方：Fork \(\rightarrow\) make} 

  \begin{block}{準備}
    \begin{enumerate}
    \item
      このレポジトリを \emph{Fork} する
      \begin{itemize}
        \item 右上の Fork \(\rightarrow\) フォーク先 として自分を選択
      \end{itemize}
    \item
      ローカルでビルドする場合，\\
      通常の \LaTeX 環境に加えて \emph{latexmk をインストールする}
      \begin{itemize}
        \butitem そうでない場合は docker をインストールする
      \end{itemize}
    \end{enumerate}
  \end{block}

  \begin{block}{pdf のビルド}
  latexmk を使う場合は\ \raise1pt\hbox{\fbox{\lstinline[language = sh]|make|}}
  \begin{itemize}
    \butitem
    docker を使う場合は\ \raise1pt\hbox{\fbox{\lstinline[language = sh]|make ondocker|}}
  \end{itemize}
  \end{block}
\end{frame}

\section{はじめてのスライド}

\begin{frame}[fragile]
  \frametitle{はじめてのスライド：簡単なサンプル}
  \begin{columns}
    \begin{column}{0.55\textwidth}
      \lstinputlisting[
        language = tex,
        linewidth = \textwidth
      ]{./fig/sample-02.tex}
    \end{column}
    \begin{column}{0.45\textwidth}
      \begin{itemize}
        \item
        全体の文字の大きさを\\
        指定できる{\footnotesize （現在 10 pt）}
        \item
        コメントは \alert{\% ではじめる}
        \item
        スペースをあけたいときは\\
        「\textbackslash\textvisiblespace」とする
        \item
        Beamer では，\\
        一枚一枚のスライドを\\
        「\Emph{フレーム}」という
      \end{itemize}
    \end{column}
  \end{columns}

\end{frame}

\begin{frame}[fragile]
  \frametitle{はじめてのスライド：箇条書き・リスト}
  \begin{columns}
    \begin{column}{0.51\textwidth}
      \lstinputlisting[
        language = tex,
        linewidth = \textwidth
      ]{./fig/sample-03.tex}
    \end{column}
    \begin{column}{0.49\textwidth}
      箇条書きには itemize を使う
  
      \begin{itemize}
      \item これは 1 番目
      \item これは 2 番目
      \item これは 3 番目
      \end{itemize}
    \end{column}
  \end{columns}

  \begin{columns}
    \begin{column}{0.51\textwidth}
      \lstinputlisting[
        language = tex,
        linewidth = \textwidth
      ]{./fig/sample-04.tex}
    \end{column}
    \begin{column}{0.49\textwidth}
      順序があるときは enumerate を使う
  
      \begin{enumerate}
      \item まずあれやって
      \item 次にこれやる
      \end{enumerate}
    \end{column}
  \end{columns}

\end{frame}

\begin{frame}[fragile]
  \frametitle{はじめてのスライド：箇条書きを{\normalsize ちょっと}良くするコツ}
  \begin{columns}
    \begin{column}{0.54\textwidth}
      \lstinputlisting[
        language = tex,
        linewidth = \textwidth
      ]{./fig/sample-05.tex}
    \end{column}
    \begin{column}{0.46\textwidth}
      \begin{itemize}
        \ngitem
        全部箇条書きだと，\\
        項目間の関係がわからない
        \thusitem 記号をうまく使うと良いかも
      \end{itemize}
    \end{column}
  \end{columns}

  \begin{mybox}{\texttt{\textbackslash item} マクロ一覧%
        {\footnotesize （勝手に定義した）}
    }
    \begin{itemize}
      \okitem \verb|\okitem|: OK !
      \ngitem \verb|\ngitem|: NG
      \thusitem \verb|\thusitem|: Thus, \dots
      \butitem \verb|\butitem|: But, \dots
      \egitem \verb|\egitem|: For example, \dots
    \end{itemize}
  \end{mybox}
\end{frame}

\begin{frame}[fragile]
  \frametitle{はじめてのスライド：強調}
  \begin{enumerate}
  \item
    \emph{強調} するときは \verb|\emph{強調}|%
    \footnote{\insertfootnotemark
    ぼくが勝手に定義を上書きした（太くなるようにしてみた）
    }
  \item
    \Emph{すごい強調} するときは \verb|\Emph{すごい強調}|%
    \footnote{\insertfootnotemark
    ぼくが勝手に定義した
    }
  \item \EMPH{めちゃめちゃすごい強調} するときは \verb|\EMPH{めちゃめちゃすごい強調}|%
    \footnote{\insertfootnotemark
    ぼくが勝手に定義した
    }
  \end{enumerate}
  
\end{frame}


\begin{frame}[fragile]
  \frametitle{はじめてのスライド：ブロック・定義・例}
  \begin{columns}
    \begin{column}{0.52\textwidth}
      \lstinputlisting[
        language = tex
      ]{./fig/sample-06.tex}
    \end{column}
    \begin{column}{0.48\textwidth}
      \begin{block}{ブロックのタイトル}
        本来ブロックは背景とは違う色で塗りつぶされるが，
        ごちゃごちゃして見えるのでやめてみた
      \end{block}
    \end{column}
  \end{columns}
  \begin{columns}
    \begin{column}{0.52\textwidth}
      \lstinputlisting[
        language = tex
      ]{./fig/sample-07.tex}
    \end{column}
    \begin{column}{0.48\textwidth}
      \begin{definition}
        1と自分自身しか約数を持たない，\\
        1 より大きい自然数を \emph{素数} という
      \end{definition}

      \begin{example}   
        \begin{itemize}
          \okitem 2 は素数．
          \okitem 3 も素数．
          \ngitem 4 は素数ではない．
        \end{itemize}
      \end{example}
    \end{column}
  \end{columns}
\end{frame}

\begin{frame}[fragile]{はじめてのスライド：ソースコード}
  \lstinputlisting[
    language = tex,
    linewidth = 0.7\textwidth
  ]{./fig/sample-10.tex}
  \begin{itemize}
    \item
    frame のオプションに \Emph{fragile を渡す必要がある}
    \item
    ソースファイルは，Makefile があるところからの相対パスで指定する\\
    （この例では，\raise1pt\hbox{\fbox{\lstinline[language=sh]|./fig/sample-10.tex|}}）
    \item
    オプションは色々ある
    \begin{itemize}
      \egitem language, linewidth, $\dots$
    \end{itemize}
  \end{itemize}

\end{frame}

\begin{frame}[fragile]
  \frametitle{はじめてのスライド：画像の挿入}
  \begin{columns}
    \begin{column}{0.67\textwidth}
      \lstinputlisting[
        language = tex,
        linewidth = \textwidth
      ]{./fig/sample-11.tex}
    \end{column}
    \begin{column}{0.33\textwidth}
      \includegraphics[width = 3cm]{./fig/logo}
    \end{column}
  \end{columns}
    ファイルは，Makefile があるところからの相対パスで指定する\\
    （この例では，\raise1pt\hbox{\fbox{\lstinline[language=sh]|./fig/logo.jpg|}}）  
\end{frame}

\section{ちょっとカスタマイズ}

{
\definecolor{maincolor}{RGB}{210, 100, 130}
\definecolor{accentcolor}{RGB}{0, 170, 210}

\begin{frame}{ちょっとカスタマイズ：色のテーマを変更}
  \begin{block}{色の変更}
    \Emph{maincolor}，\EMPH{accentcolor} を設定すると，\\
    スライド全体の色を変えられます
    \lstinputlisting[
      language = tex,
      linewidth = \textwidth
    ]{./fig/sample-13.tex}
    \begin{itemize}
      \ngitem \textcolor{blue}{原色} そのままとかじゃなくて
      \okitem かっこいい色を探してみてください
    \end{itemize}
  \end{block}

\end{frame}
}

{
\setbeamertemplate{headline}{}
\setbeamercolor{footline}{fg = white}

\begin{frame}{ちょっとカスタマイズ：ヘッダとプログレスメータ}
\begin{block}{ヘッダ}
    スライド上部の \emph{ヘッダ} は，
    \lstinputlisting[
      language = tex,
      linewidth = \textwidth
    ]{./fig/sample-14.tex}
    と書いてやると消えます
    \begin{itemize}
    \okitem
      スライドが始まる前に書いておけば\\
      無駄なスペースができることもないです
    \end{itemize}
\end{block}

\begin{block}{プログレスメータ}
    スライド下部の \emph{プログレスメータ} は，
    \lstinputlisting[
      language = tex,
      linewidth = \textwidth
    ]{./fig/sample-12.tex}
    と書いてやると（背景と同じ色になって）消えます
  \end{block}

\end{frame}
}

